% Seitenrand: Innen mind. 25mm
% Außenrand: Korrekturrand 50mm,
%   Randnotizen können in den 50mm-Rand hineinragen

\subsubsection{Papierformat}


\subsubsection{Ränder}

Die Seitenränder sollten gemäß DIN~5008 mind. \unit[2,5]{cm} betragen. Bei
Arbeiten, die auf der Außenseite einen Korrekturrand erfordern, sollte der
Korrekturrand \unit[5]{cm} breit sein.

Die Einstellung des Satzspiegels erfolgt mit dem Paket \texttt{geometry}.

\newpage
\setlayoutscale{0.29}
\printheadingsfalse
\printparametersfalse
%\drawdimensionstrue
\currentpage
\begin{minipage}{0.5\linewidth}
\oddpagelayoutfalse
\pagedesign
\end{minipage}%
\begin{minipage}{0.5\linewidth}
\oddpagelayouttrue
\pagedesign
\end{minipage}


\subsubsection{Duplexdruck}

Sofern nicht explizit anders gefordert, sollten Arbeiten beidseitig gedruckt
werden, um Papier zu sparen. Hierzu ist die Klassenoption \texttt{twoside} zu
setzen.

\todo[inline]{Code-Beispiel}

Für kürzere Haus- oder Seminararbeiten ist im Zweifel auch einseitiger Druck
statthaft.

In umfangereicheren Arbeiten (ab Bachelorarbeit aufwärts), sollten neue
Kapitel stets auf einer rechten Seiten beginnen. Die Klasse \texttt{book} wird
vor Beginn eines neuen Kapitels im Bedarfsfall automatisch eine Leerseite
einfügen. Bei der Klasse \texttt{report} ist hierfür die Option
\texttt{openright} notwendig.

\todo[inline]{Code-Beispiel}




\subsubsection{Bindekorrektur}

Für ein ordentliches Erscheinugsbild werden Studienabschlussarbeiten in der
Regel in einem Hardcover-Einband eingereicht. Zur Anwendung kommt meist eine
Kamm- oder Klebebindung. Da ein Teil des Innenrandes von der Bindung verdeckt
wird, muss der Satzspiegel entsprechend ein klein wenig enger gestaltet werden,
um ein optisch stimmiges Erscheinungsbild zu erzielen. Diese Bindekorrektur
bestimmt sich wie folgt \dots

\todo[inline]{Formeln}


\todo[inline]{Code-Beispiel, Formeln}

Bei Ringbindungen muss keine Bindekorrektur erfolgen.

% DTK 2004/1

% Bindekorrektur bei Klebe- und Kammbindung: 1mm pro 10 Blatt bei 80g/m^2
